\section{Деревья решений (полные и сокращенные). Синтаксические деревья. Бинарная диаграмма решений.}

\subsection*{Синтаксические деревья}

\paragraph{Определение}
\textbf{Синтаксическое дерево} — это иерархическое представление структуры \textbf{формулы}. 
\begin{itemize}
    \item Корнем является главная (внешняя) операция формулы.
    \item Узлами — промежуточные операции.
    \item Листьями — переменные или константы.
\end{itemize}
Это дерево отражает порядок вычислений (интерпретации) по алгоритму Eval: сначала вычисляются значения детей, затем — значение родителя.

\subsection*{Деревья решений (семантические деревья)}

\paragraph{Полное дерево решений}
Таблицу истинности функции от $n$ переменных можно представить в виде полного бинарного дерева высоты $n+1$.
\begin{itemize}
    \item \textbf{Ярусы:} соответствуют переменным $x_1, \dots, x_n$ в установленном порядке.
    \item \textbf{Дуги:} левая дуга соответствует значению 0, правая — 1.
    \item \textbf{Листья:} хранят значения функции $f$ на кортеже, соответствующем пути из корня.
\end{itemize}

\paragraph{Сокращенное дерево решений}
Если у некоторого узла все листья в его поддереве имеют одно и то же значение (все 0 или все 1), то всё поддерево можно заменить этим значением. Это значительно уменьшает объём хранимых данных.

\subsection*{Бинарная диаграмма решений (BDD)}

\paragraph{Определение}
\textbf{Бинарная диаграмма решений} получается из дерева решений путём отказа от древовидности связей (допускается несколько входящих дуг в один узел). Построение диаграммы включает три преобразования:
\begin{enumerate}
    \item \textbf{Отождествление листьев:} Все листовые узлы, содержащие 0, сливаются в один узел «0». Аналогично — для «1».
    \item \textbf{Слияние изоморфных поддиаграмм:} Если два узла одного яруса ведут в одни и те же узлы при одинаковых значениях переменной, они заменяются одним экземпляром.
    \item \textbf{Исключение избыточных узлов:} Если обе дуги из узла (0 и 1) ведут в один и тот же следующий узел, этот узел удаляется, а входящие в него дуги перенаправляются напрямую к следующему узлу.
\end{enumerate}

\paragraph{Зависимость от порядка переменных}
Результат построения диаграммы существенно зависит от порядка переменных. Правильный выбор порядка (например, вынесение наиболее влияющей переменной в корень) может сделать диаграмму экспоненциально компактнее.
\textit{Пример Новикова:} Функция, которая при одном порядке требует сложной диаграммы, при другом (например, $y, x, z$) может превратиться в простое условие `if y then 1 else !x`.

\subsection*{Алгоритм интерпретации}

Вычисление значения функции по дереву или диаграмме (Алгоритм 3.5) выполняется проходом от корня к листу:
\begin{itemize}
    \item На каждом шаге смотрим на значение текущей переменной $x_i$.
    \item Если $x_i = 0$, переходим по левой дуге, если $x_i = 1$ — по правой.
    \item Когда достигнут листовой узел, возвращается его значение.
\end{itemize}

\paragraph{Преимущество}
Представление функции в виде BDD часто позволяет хранить меньше информации и производить вычисления быстрее, чем при использовании традиционных таблиц или формул.